\documentclass[11pt,a4paper]{article}
\usepackage[margin=1in]{geometry}
\usepackage{amsmath,amssymb,amsfonts}
\usepackage{mathtools}
\usepackage{lmodern}
\usepackage[T1]{fontenc}
\usepackage[utf8]{inputenc}
\usepackage{textcomp}
\usepackage{microtype}
\usepackage{parskip}
\usepackage{enumitem}
\usepackage{array}
\usepackage{booktabs}
\usepackage{hyperref}
\hypersetup{hidelinks,
  pdftitle={Paper 090 -- The V5 Cross-Chain Correlation Kernel},
  pdfauthor={Allen Proxmire}}
\setlength{\parindent}{0pt}
\setlength{\parskip}{6pt plus 2pt minus 1pt}

\title{\vspace{-1cm}\Large\bfseries Paper 090 --- The V5 Cross-Chain Correlation Kernel}
\author{Allen Proxmire}
\date{May 2026}

\begin{document}
\maketitle

\section*{Abstract}

The Event Density (ED) substrate is a 13-primitive generative system whose downstream content includes a hierarchy of kernel rule-types. The first such kernel, \textbf{V1}, is the substrate's single-chain vacuum-response kernel (finite-width per Theorem N1, Paper \#18; retarded-support per Theorem T18, Paper\_093). This paper introduces and characterizes the second kernel rule-type: \textbf{V5}, the \textbf{cross-chain correlation kernel}. Given the postulated primitives P02, P04, P05, P07, P09, P10, P11, together with V1's finite-width retarded structure, V5 is characterized as the substrate kernel rule-type that mediates \emph{correlation between distinct chains} with finite memory time $\tau_{V5}$, forward-causal support, and $U(1)$-gauge-compatible transport.

\textbf{Wedge claim.} V5 provides a \textbf{structural unification of formalism across regimes with different correlation times}: the same substrate kernel rule-type --- with the same structural features (finite memory, retarded support, gauge-compatibility, Lorentz covariance) --- supplies a unified substrate mechanism for non-Markovian behavior across three physical regimes: (i) Maxwell viscoelastic stress-relaxation in polymer melts and non-Newtonian fluids; (ii) the substrate-level cutoff producing the Hawking spectrum's trans-Planckian regularization at near-horizon scales; (iii) cross-chain entanglement-bandwidth modulation that controls Page-curve behavior and ER=EPR-class straddling at horizons.

\textbf{The regime-specific memory times $\tau_{V5}^\mathrm{soft}$, $\tau_{V5}^\mathrm{BH}$, $\tau_{V5}^\mathrm{ent}$ are inherited parameters, not predicted by the substrate.} Specifically, $\tau_{V5}^\mathrm{soft}$ is \emph{identified with} the empirically-measured Maxwell relaxation time, and $\tau_{V5}^\mathrm{BH}$ is \emph{chosen to match} the substrate scale $\ell_P/c$. These are operational identifications/matchings, not substrate-level predictions of the numerical values.

The Diffusion Coarse-Graining Theorem (DCGT) bridges substrate-level V5 to continuum constitutive laws in each regime. The paper is conditional: given the seven listed primitives + V1 inheritance, V5 is characterized; the primitives themselves are postulated. The cross-regime structural unification of formalism is the empirical case.

\section{Introduction}

\subsection{What V5 is}

V5 is a \textbf{substrate-level kernel rule-type} distinct from V1 (single-chain vacuum kernel) and distinct from matter rule-types (chain-carrying participation measures). Where V1 mediates a single chain's vacuum response, V5 mediates \textbf{correlations between distinct chains} through the substrate. Operationally, V5 is the substrate-level mechanism by which one chain's participation content can influence another chain's participation content across substrate-graph distance, subject to finite memory time, forward-causal direction, and the gauge structure of P05--P09.

V5 is \emph{not} a free-standing postulate of ED. It is a rule-type the substrate carries given the 13 postulated primitives + V1 inheritance: P10 (multiple structurally distinct rule-types) opens room for kernel rule-types beyond V1; P02 + P04 + P05 + P07 + P09 supply the structural mechanism for inter-chain correlation transport; P11 + V1's retardation supply the causal-direction constraint.

\subsection{Why cross-chain kernels matter}

Standard physics treats correlations between distinct subsystems via field-theoretic propagators. In ED's substrate ontology, there are no fundamental fields (ED-I-06); chains and channels are the primitive ontological objects, and correlations between distinct chains are mediated by \emph{substrate-level kernel rule-types} rather than by continuum fields. V5 is the substrate-level analog of the cross-chain ``field propagator'' operating on the discrete participation graph rather than on a smooth continuum.

The interest of V5 is its \textbf{structural unification of formalism across regimes with different correlation times}: the same substrate rule-type (with the same structural features) provides a unified substrate mechanism for non-Markovian behavior across three physical regimes:

\begin{itemize}[noitemsep]
\item \textbf{Soft-matter physics:} Maxwell viscoelastic stress relaxation in polymer melts, complex fluids, non-Newtonian rheology.
\item \textbf{Black-hole physics:} the substrate-level cutoff at near-horizon scales that regularizes the Hawking spectrum's trans-Planckian extrapolation.
\item \textbf{Quantum entanglement:} cross-chain bandwidth that controls entanglement-entropy growth, Page-curve behavior, ER=EPR-class straddling at horizons.
\end{itemize}

The three regimes have \textbf{different correlation times} ($\tau_{V5}^\mathrm{soft}$, $\tau_{V5}^\mathrm{BH}$, $\tau_{V5}^\mathrm{ent}$, all regime-specific inherited parameters), but the \textbf{substrate-level kernel structure is the same} in all three: finite memory, retarded support, gauge-compatibility, cross-chain operational form. This structural unification of formalism is the wedge claim of V5.

\subsection{V5 in the kernel hierarchy}

\begin{itemize}[noitemsep]
\item \textbf{V1} (Papers \#18 / Paper\_093): single-chain vacuum response. Finite-width, retarded, Lorentz-covariant.
\item \textbf{V5} (this paper): cross-chain correlation. Finite-memory $\tau_{V5}$, retarded, gauge-compatible.
\item \textbf{N1-E, N2-E, N3-D} (future Wave-2 papers): memory-kernel cascade.
\item \textbf{Lindblad limit} (future Wave-2 paper): Markovian coarse-grained limit of V1 + V5.
\end{itemize}

V1 and V5 are the two foundational kernels; the cascade and Lindblad limit are downstream of their joint action.

\subsection{What this paper does}

\S2 lists primitive inputs. \S3 defines V5. \S4 establishes structural properties. \S5 develops three regimes (soft-matter, Hawking, entanglement). \S6 develops DCGT. \S7 names what the paper does \emph{not} claim. \S8 lists falsification criteria. \S9 is the position statement.

\section{Primitive Inputs (postulated, not derived in this paper)}

\begin{itemize}[noitemsep]
\item \textbf{P02 (Participation as primitive relation).}
\item \textbf{P04 (Bandwidth as non-negative additive scalar).}
\item \textbf{P05 (Polarity-transport along edges).}
\item \textbf{P07 (Channel structure as ontological primitive).}
\item \textbf{P09 ($U(1)$-valued polarity).}
\item \textbf{P10 (Rule-type primitive).}
\item \textbf{P11 (Commitment irreversibility).}
\end{itemize}

\textbf{V1 inheritance:}

\begin{itemize}[noitemsep]
\item V1 finite-width (Theorem N1, Paper \#18).
\item V1 retarded support (Theorem T18, Paper\_093).
\end{itemize}

\textbf{Upstream paper dependencies:} Papers \#1, \#3, \#5, \#18, \#093.

\textbf{No primitive forcing is invoked.}

\section{Definition of the V5 Kernel}

\subsection{Participation-based definition}

For two chains $C_A$, $C_B$ at distinct loci, V5 is:
\[
K_{V5}(u_A, t_A; u_B, t_B) = \theta(t_A - t_B)\, F_{V5}\!\left(\frac{\sigma(u_A, u_B; t_A, t_B)}{\ell_{V5}^{\,2}},\, \frac{t_A - t_B}{\tau_{V5}}\right)
\]
where:

\begin{itemize}[noitemsep]
\item $\theta(t_A - t_B)$: Heaviside step enforcing retarded support (forward-causal from V1 via P11).
\item $\sigma(u_A, u_B; t_A, t_B)$: substrate-level Synge function analog.
\item $\ell_{V5}$: V5-characteristic spatial scale (regime-dependent; inherited from $\ell_\mathrm{ED}$ at substrate scale, renormalized in coarse-grained regimes via DCGT).
\item $\tau_{V5}$: V5-characteristic memory time. \textbf{Regime-specific; inherited, not predicted (\S3.3 below).}
\item $F_{V5}$: V5 envelope function --- bounded, decaying. The \emph{form} (bounded, finite-width, retarded) is fixed by V5's substrate-rule-type identity; the \emph{specific functional shape} (exponential, power-law, multi-scale) is value-layer content.
\end{itemize}

The cross-chain correlation acts:
\[
\langle X^A(u_A, t_A) \cdot Y^B(u_B, t_B)\rangle_{V5} = \int K_{V5}(u_A, t_A; u_B, t_B)\,\mathcal{Q}_{X,Y}^{AB}\,d\mu(\mathrm{intermediate}).
\]

\subsection{Cross-chain correlation structure}

V5 acts on \textbf{pairs of chains}, distinguishing it from V1 (single-chain operator). The two-chain character is essential; without V5, the substrate would have only V1 + matter rule-types, and cross-chain correlations would require either repeated V1 cascade (Markovian only) or substrate-level fundamental fields (forbidden by ED-I-06).

V5 provides a unified substrate mechanism for non-Markovian cross-chain correlation transport without invoking continuum fields. \textbf{The substrate mechanism is unified across regimes}; the specific $\tau_{V5}$ in each regime is inherited from value-layer empirical content.

\subsection{Memory time $\tau_{V5}$ --- regime-specific, inherited, not predicted}

The memory-time parameter $\tau_{V5}$ controls cross-chain correlation decay:
\[
F_{V5}\!\left(\sigma/\ell_{V5}^{\,2}, t/\tau_{V5}\right) \xrightarrow{t/\tau_{V5} \to \infty} 0.
\]

\textbf{Structural content (substrate-derived):}

\begin{itemize}[noitemsep]
\item The \emph{existence} of a finite memory time is structural (forced by V1's finite-width inheritance + UV-finiteness of substrate primitives).
\item The \emph{form} of the V5 kernel (finite-memory, retarded, gauge-compatible, Lorentz-covariant) is structural.
\end{itemize}

\textbf{Value-layer content (inherited, not substrate-derived):}

\begin{itemize}[noitemsep]
\item The \emph{specific numerical value} of $\tau_{V5}$ in any given physical system is \textbf{inherited from value-layer empirical content}, not predicted by the substrate.
\item The \emph{specific decay profile} (exponential, power-law, multi-scale) within the structural class is inherited.
\end{itemize}

In the regime-specific evaluations of \S5, three regime-specific values appear: $\tau_{V5}^\mathrm{soft}$, $\tau_{V5}^\mathrm{BH}$, $\tau_{V5}^\mathrm{ent}$. \textbf{All three are inherited parameters, not substrate predictions.} Specifically:

\begin{itemize}[noitemsep]
\item \textbf{Soft-matter regime:} $\tau_{V5}^\mathrm{soft}$ is \textbf{identified with} the empirically-measured Maxwell relaxation time of polymer melts and non-Newtonian fluids ($\sim 10^{-3}$ s for typical samples). The substrate framework provides the operational mechanism (V5 cross-chain correlation); the numerical value is inherited from rheometry.
\item \textbf{Black-hole / Hawking regime:} $\tau_{V5}^\mathrm{BH}$ is \textbf{chosen to match the substrate scale} $\ell_P/c \sim 10^{-44}$ s. This is a substrate-scale matching, not a substrate-level derivation: $\ell_P$ itself is empirically identified with $\ell_\mathrm{ED}$ via Newton-recovery (Paper\_027), so $\tau_{V5}^\mathrm{BH} = \ell_P/c$ matches the substrate scale by construction. The Hawking-cutoff result of \S5.2 is conditional on this matching.
\item \textbf{Entanglement regime:} $\tau_{V5}^\mathrm{ent}$ is system-specific, inherited from the specific quantum-information context.
\end{itemize}

The substrate's structural prediction is that V5's memory time has a \textbf{consistent operational definition across regimes} (same substrate kernel rule-type, same structural features), but the \textbf{numerical value in any specific system is inherited from value-layer empirical content}.

\subsection{Relation to V1}

V5 inherits from V1:

\begin{enumerate}[noitemsep]
\item \textbf{Finite-width support (V1 Theorem N1):} $F_{V5}$ bounded, decaying outside V5-characteristic scale. V5 has no $\delta$-function limit and no infinite-width limit.
\item \textbf{Retarded support (V1 Theorem T18 / Paper\_093):} $K_{V5}$ has $\theta(t_A - t_B)$ support. Forced by the same P11 commitment-irreversibility argument that fixes V1's retardation. Advanced / symmetric V5 are excluded by the same substrate-operational compositional-closure arguments (Paper\_093 \S5.3, \S6.4).
\end{enumerate}

\subsection{V5 is a substrate rule-type}

By P10, V5 occupies the cross-chain-correlation slot in the kernel-rule-type category. V5 is not a matter rule-type, not a gauge rule-type --- it is its own structural category alongside V1 and the future cascade kernels.

\section{Structural Properties}

\subsection{Linearity}

V5 acts linearly on cross-chain participation content.

\subsection{Boundedness}

V5 produces bounded cross-chain correlations.

\subsection{Gauge-compatibility via P05 + P09}

Under $U(1)$ gauge transformations, $K_{V5}$ transforms covariantly:
\[
K_{V5} \to e^{i(\alpha(u_A) - \alpha(u_B))}\,K_{V5}.
\]

\subsection{Lorentz covariance}

V5 transforms as a Lorentz scalar; envelope $F_{V5}$ depends on Lorentz-invariant separation. In curved spacetime, generalizes via Synge world-function on $g_{\mu\nu}$.

\subsection{Rule-type identity via P10}

V5 is a kernel rule-type alongside V1 in the rule-type category.

\subsection{Retarded support and the kernel-level arrow}

V5's retarded support inherits from V1 via P11 (Paper\_093 \S6.4 substrate-operational argument).

\section{Cross-Domain Manifestations of V5}

V5 provides a \textbf{unified substrate mechanism for non-Markovian behavior across regimes} with different correlation times. The three regimes have \textbf{different inherited $\tau_{V5}$ values} but share the \textbf{same substrate-level kernel formalism}.

\subsection{Soft-matter Maxwell viscoelastic relaxation}

\textbf{The phenomenon.} In polymer melts, complex fluids, and non-Newtonian rheology, stress relaxation follows the Maxwell law:
\[
\sigma(t) = \sigma_0\,e^{-t/\tau_\mathrm{Maxwell}}.
\]
In typical polymer-melt systems, $\tau_\mathrm{Maxwell}$ is in the $10^{-3}$ s range, varying with molecular weight, temperature, chemistry.

\textbf{The V5 identification.} Polymer-melt molecules are extended chain structures in the substrate sense (P02 + P07). Cross-chain correlations between distinct polymer molecules are mediated at the substrate level by V5. The substrate-level V5 envelope produces an exponential-decay form in the soft-matter regime:
\[
F_{V5}(\sigma/\ell_{V5}^{\,2}, t/\tau_{V5}^\mathrm{soft}) \sim e^{-t/\tau_{V5}^\mathrm{soft}}
\]
at large temporal separations. Through DCGT (\S6), this becomes the continuum Maxwell stress-relaxation law.

\textbf{Identification of $\tau_{V5}^\mathrm{soft}$.} \textbf{$\tau_{V5}^\mathrm{soft}$ is identified with the empirically-measured Maxwell relaxation time $\tau_\mathrm{Maxwell}$, not predicted by the substrate.} Specifically:

\begin{itemize}[noitemsep]
\item The substrate framework supplies the \emph{operational mechanism} (V5 cross-chain correlation with exponential-envelope structural form) that produces the Maxwell stress-relaxation law via DCGT.
\item The \emph{numerical value} of $\tau_{V5}^\mathrm{soft}$ in any specific polymer system is inherited from rheometry --- it is the empirically-measured Maxwell relaxation time for that system.
\item ED does \emph{not} predict the specific value of $\tau_\mathrm{Maxwell}$ from molecular structure (molecular weight, temperature, chemistry). That prediction would require additional substrate-level content connecting molecular structure to substrate-graph cross-chain correlation parameters, which is in-queue work.
\end{itemize}

\textbf{The substrate framework's structural prediction is therefore:} the \textbf{form} of stress-relaxation (exponential, Maxwell-type) in the soft-matter regime is a substrate-level structural consequence of V5's finite-memory envelope; the \textbf{value} of the relaxation time is inherited from value-layer empirical content.

\textbf{Arc-D context.} Substrate-level V5-Maxwell identification developed in Arc-D / DCGT memos.

\subsection{Hawking-spectrum cutoff via imaginary-time periodicity}

\textbf{The phenomenon.} Hawking radiation is thermal flux at $T_H = \kappa/(2\pi)$. Standard derivation requires extrapolation of matter-field modes to $\omega \to \infty$ --- the trans-Planckian problem.

\textbf{The V5 identification.} Near the horizon, V5 acts as the substrate-level cross-chain correlation kernel between chains carrying Hawking-radiation excitations and substrate near-horizon vacuum structure. V5 contributes the substrate-level KMS condition:
\[
\langle\phi(u_A, t_A)\phi(u_B, t_B)\rangle_{V5}^\mathrm{Eucl} = \langle\phi(u_A, t_A + i\beta_H)\phi(u_B, t_B)\rangle_{V5}^\mathrm{Eucl}, \quad \beta_H = 2\pi/\kappa,
\]
producing thermal Planck distribution at $T_H = 1/\beta_H = \kappa/(2\pi)$.

The cutoff frequency that regularizes the trans-Planckian extrapolation is set by V5's substrate-level memory time:
\[
\omega_c = \frac{1}{\tau_{V5}^\mathrm{BH}} \sim \frac{c}{\ell_P}.
\]

\textbf{The identification $\tau_{V5}^\mathrm{BH} = \ell_P/c$ is chosen to match the substrate scale, not derived from substrate-level analysis.} Specifically:

\begin{itemize}[noitemsep]
\item The substrate framework requires V5's near-horizon memory time to be at the substrate scale (consistent with V5's substrate-level kernel-rule-type identity).
\item $\ell_P$ is empirically identified with $\ell_\mathrm{ED}$ via Newton-recovery (Paper\_027 \S5.3); this is a value-layer empirical identification, not a substrate-level derivation.
\item The matching $\tau_{V5}^\mathrm{BH} = \ell_P/c$ is therefore an \textbf{operational substrate-scale matching}, conditional on $\ell_\mathrm{ED} = \ell_P$.
\end{itemize}

If $\ell_\mathrm{ED}$ were identified with a different empirical scale, $\tau_{V5}^\mathrm{BH}$ would shift correspondingly. The Hawking-cutoff result is conditional on this matching.

\textbf{The substrate framework's structural content (independent of the specific numerical value):}

\begin{itemize}[noitemsep]
\item V5 provides a substrate-level kernel mechanism that \emph{structurally produces} a UV cutoff in the Hawking spectrum (from V5's finite-memory envelope).
\item The cutoff frequency is the inverse of V5's near-horizon memory time, scaled to the substrate scale.
\item The trans-Planckian extrapolation is not generated because V5-mediated cross-chain correlations cannot extrapolate to frequencies $\omega > 1/\tau_{V5}^\mathrm{BH}$.
\end{itemize}

The numerical value $\omega_c \sim c/\ell_P$ matches the physical Planck scale \textbf{because} the substrate framework identifies $\ell_\mathrm{ED}$ with $\ell_P$ --- this is empirical matching, not substrate derivation.

\textbf{Arc-Hawking context.} Substrate-level V5-Hawking-cutoff developed in Arc-Hawking memos (H-4 closure).

\subsection{Entanglement-bandwidth modulation}

\textbf{The phenomenon.} Entanglement entropy growth between subsystems is bounded by cross-chain correlation capacity. In black-hole physics, the Page curve describes entanglement-entropy growth and saturation. In condensed-matter physics, entanglement-entropy growth follows analogous patterns bounded by the Lieb-Robinson velocity.

\textbf{The V5 derivation.} Entanglement between chains $C_A$, $C_B$ is carried by V5-mediated cross-chain correlation content. Maximum entanglement-bandwidth growth rate:
\[
\frac{dS_E}{dt}\bigg|_\mathrm{max} = c_{V5} \cdot \frac{\sigma_{AB}}{\tau_{V5}^\mathrm{ent}}.
\]
V5's finite-memory structure produces entanglement-entropy saturation at a characteristic timescale related to $\tau_{V5}^\mathrm{ent}$.

In black-hole physics, Page-curve transition occurs at $t_\mathrm{Page} \sim \tau_{V5}^\mathrm{ent}(M)$; substrate-level mechanism for ER=EPR-class straddling. In condensed-matter, the Lieb-Robinson bound emerges as the DCGT continuum limit of V5.

\textbf{Empirical identification:} $\tau_{V5}^\mathrm{ent}$ is system-specific, \textbf{inherited from value-layer empirical content}, not predicted by the substrate.

\textbf{Arc-E context.} Arc-E memos (E-5 closure).

\subsection{The structural unification (summary)}

V5 provides a unified substrate kernel rule-type with consistent structural features (finite memory, retarded, gauge-compatible, Lorentz-covariant) across three regimes with different inherited correlation times:

\begin{center}
\begin{tabular}{p{0.30\textwidth}p{0.30\textwidth}p{0.32\textwidth}}
\toprule
\textbf{Regime} & $\tau_{V5}$ \textbf{status} & \textbf{Substrate role} \\
\midrule
Soft-matter polymer melts & Identified with empirical $\tau_\mathrm{Maxwell}$ & V5 mechanism for Maxwell viscoelastic relaxation \\
Near-horizon black hole & Chosen to match substrate scale $\ell_P/c$ & V5 mechanism for Hawking spectrum cutoff \\
Entanglement (system-dependent) & Inherited from system & V5 mechanism for entanglement-bandwidth modulation \\
\bottomrule
\end{tabular}
\end{center}

\textbf{The structural unification of formalism is the wedge claim.} What V5 supplies across these regimes is the \textbf{same substrate kernel structure} (the same primitive content, the same kernel-rule-type identity, the same set of structural properties --- finite memory, retarded support, gauge-compatibility, cross-chain operational form), with \textbf{regime-specific inherited correlation times}.

V5 is \textbf{not} claimed to \emph{necessarily} underlie all non-Markovian behavior. Standard frameworks (polymer-chain reptation, mode-coupling theory, Lieb-Robinson bounds derived from microscopic Hamiltonians) provide alternative descriptions in their respective regimes. V5 provides a \textbf{unified substrate mechanism} for these phenomena --- a single substrate-level kernel structure that produces these effects via DCGT. Whether V5 is the \emph{only} possible substrate-level mechanism, or whether alternative substrate ontologies could supply equally consistent mechanisms, is open.

The empirical case for V5 (and for the 13-primitive ED ontology) is the structural unification of formalism: one substrate kernel rule-type doing structurally identical jobs in three regimes that differ in correlation time. This is the wedge content.

\section{Diffusion Coarse-Graining (DCGT) and Continuum Constitutive Laws}

DCGT (Arc-D closure) bridges substrate-level V5 + V1 + matter-rule-type content to continuum constitutive laws in the hydrodynamic-window scale separation $\ell_\mathrm{ED} \ll R_\mathrm{cg} \ll L_\mathrm{flow}$.

\subsection{V5 $\to$ Maxwell relaxation in soft matter}

DCGT coarse-graining of V5 cross-chain correlations in soft-matter produces:
\[
G(t) = G_0\,e^{-t/\tau_\mathrm{Maxwell}}
\]
with $\tau_\mathrm{Maxwell}$ identified with the substrate-level $\tau_{V5}^\mathrm{soft}$ rescaled by the DCGT coarse-graining kernel. Krieger-Dougherty viscosity and Cross / Carreau / Bird-Carreau constitutive laws emerge as DCGT continuum limits.

\subsection{V5 $\to$ propagator structure in QFT}

DCGT coarse-graining of V5 in the QFT regime produces retarded / Feynman / advanced propagator hierarchy. V5's substrate-level UV cutoff at $\omega \sim 1/\tau_{V5}$ becomes standard QFT regulators (Pauli-Villars, dimensional, lattice) in the continuum limit. In near-horizon regime, V5's $\omega_c \sim c/\ell_P$ (matched to substrate scale, \S5.2) provides the trans-Planckian cutoff.

\subsection{V5 $\to$ Lieb-Robinson bound in condensed matter}

DCGT coarse-graining in condensed-matter:
\[
v_{LR} = c_{V5}\cdot\frac{\ell_{V5}}{\tau_{V5}^\mathrm{ent}}.
\]

\subsection{V5 $\to$ curved-spacetime entanglement at horizons}

V5's curved-spacetime extension produces entanglement structure across horizons that ER=EPR-class arguments require: cross-chain correlations straddle the horizon, with committed structure (P11) interior-confined and uncommitted-entanglement amplitude straddling.

\subsection{DCGT structural summary}

DCGT applies to V5 wherever hydrodynamic-window scale separation holds. In each regime, the same substrate primitive produces a regime-specific continuum constitutive law via DCGT. The structural unification of formalism is preserved across these continuum limits.

\section{What This Paper Does NOT Claim}

\subsection{The V5 memory time $\tau_{V5}$ is not predicted by the substrate}

\textbf{Specific numerical values of $\tau_{V5}^\mathrm{soft}$, $\tau_{V5}^\mathrm{BH}$, $\tau_{V5}^\mathrm{ent}$ are inherited from value-layer empirical content, not predicted by the substrate}:

\begin{itemize}[noitemsep]
\item $\tau_{V5}^\mathrm{soft}$ is identified with empirically-measured Maxwell relaxation time.
\item $\tau_{V5}^\mathrm{BH} = \ell_P/c$ is chosen to match the substrate scale; $\ell_P$ itself is empirically identified with $\ell_\mathrm{ED}$ via Newton-recovery.
\item $\tau_{V5}^\mathrm{ent}$ is system-specific, inherited.
\end{itemize}

The substrate's structural content is the \textbf{kernel form + structural features} (finite memory, retarded support, gauge-compatibility), not the numerical values.

\subsection{The V5 envelope shape is not derived}

The specific functional shape of $F_{V5}$ within the admissible class is value-layer content.

\subsection{No coupling constants or masses are derived}

V5 does not derive specific material constants, BH parameters, or entanglement-coupling coefficients. V5 produces the \emph{structural form}; numerical content is value-layer inherited.

\subsection{No uniqueness claim beyond the 13 primitives}

V5's characterization is conditional on the postulated primitives + V1 inheritance. Other substrate ontologies could allow different cross-chain kernel structures.

\subsection{No forcing of primitives}

The seven primitives are postulated; not derived from a deeper layer.

\subsection{No claim of derivation from nothing}

Conditional on postulated primitives + V1 inheritance + upstream papers + DCGT.

\subsection{No claim of superiority over standard cross-correlation physics}

Standard QFT propagators, soft-matter Maxwell viscoelasticity, Lieb-Robinson bound are empirically successful. V5 does not ``replace'' them --- V5 produces them via DCGT and provides a \textbf{unified substrate mechanism for non-Markovian behavior across regimes}. The regime-specific continuum laws are inherited and DCGT-derived from V5.

\subsection{V5 is not claimed as the only possible substrate-level mechanism}

V5 provides a unified substrate mechanism for non-Markovian behavior across regimes --- within the 13-primitive ED ontology. \textbf{V5 is not claimed as necessary for non-Markovian dynamics in some absolute sense.} Standard frameworks (polymer reptation, mode-coupling theory, microscopic Lieb-Robinson derivations) supply alternative descriptions in their respective regimes; V5's distinctive content is the \emph{substrate-level unification of formalism} across these regimes within ED's ontology.

\section{Falsification Criteria}

\subsection{Failure of structural unification of formalism}

\textbf{Discovery that soft-matter Maxwell relaxation, Hawking spectrum cutoff, and entanglement-bandwidth modulation cannot be described by a single substrate-level kernel formalism with consistent structural features (finite memory, retarded support, gauge-compatibility, cross-chain operational form) --- even after accounting for regime-specific $\tau_{V5}$ --- would falsify V5's structural unification claim.}

The falsification would take the form of demonstrating that the substrate-level structural features differ across the three regimes in a way that cannot be reconciled with a single kernel rule-type. For example:

\begin{itemize}[noitemsep]
\item Soft-matter Maxwell exhibits non-finite-memory behavior at substrate scale that contradicts V5's finite-memory structure.
\item Hawking-cutoff structure exhibits non-gauge-compatible behavior contradicting V5's $U(1)$-covariance.
\item Entanglement-bandwidth modulation exhibits non-retarded behavior contradicting V5's forward-causal support.
\end{itemize}

The current state of empirical content is \textbf{consistent with single-substrate-kernel structural-formalism unification} across the three regimes. Tighter cross-regime empirical comparison could either confirm or falsify.

\subsection{Observation of independent structural-formalism mechanisms for the three regimes}

Specific empirical observations that would falsify V5's structural unification:

\begin{itemize}[noitemsep]
\item Substrate-level structural features (memory, support, gauge, Lorentz) differing between soft-matter and entanglement contexts.
\item Hawking spectrum cutoff exhibiting non-V5-compatible structural features.
\item Entanglement-bandwidth modulation exhibiting non-V5-compatible structural features.
\end{itemize}

\subsection{Observation of advanced V5 support}

Discovery of cross-chain correlations carrying advanced (backward-causal) support at the substrate level would falsify V5's retarded character (\S4.6 + Paper\_093).

\subsection{Breakdown of structural identifications}

\begin{itemize}[noitemsep]
\item \textbf{Soft-matter:} discovery that Maxwell-type stress relaxation has structurally non-V5 substrate origin would falsify the V5 identification of \S5.1.
\item \textbf{Hawking:} discovery that the trans-Planckian cutoff is at a frequency inconsistent with $\omega_c = 1/\tau_{V5}^\mathrm{BH}$ (at the substrate-scale matching) would falsify the V5 mechanism of \S5.2.
\end{itemize}

\subsection{Breakdown of gauge-compatibility}

Cross-chain correlations between gauge-coupled chains failing to transform covariantly under $U(1)$ gauge transformations would falsify \S4.3.

\subsection{Substrate-level kernel without V5}

\textbf{Discovery that the substrate-level kernel hierarchy can be exhausted by V1 alone (without V5) --- i.e., that all observed cross-chain correlation phenomena can be derived from V1 composed across intermediate chains via Markovian cascade --- would falsify V5's structural role within the substrate framework}, although it would not directly falsify the empirical phenomena V5 describes (which would be reattributed to V1-only cascade in such a counterfactual).

Current state: empirical non-Markovian cross-chain correlation structure is structurally inconsistent with single-V1 cascade composition (V5 supplies a unified substrate mechanism for non-Markovian behavior that V1-only cascade cannot supply within ED's ontology). Falsification would require demonstrating that empirical non-Markovian correlations \emph{can} be reduced to V1-only cascades.

\section{Position Statement}

V5 is a \textbf{substrate kernel rule-type} in the 13-Primitive Generative Substrate Ontology. V5's identity is \emph{conditional} on the postulated primitives P02 + P04 + P05 + P07 + P09 + P10 + P11 + V1 inheritance + upstream Wave-1 papers + DCGT.

\textbf{V5's wedge content} is the \textbf{structural unification of formalism across regimes with different correlation times}. V5 provides a unified substrate kernel rule-type --- with consistent structural features (finite memory, retarded support, gauge-compatibility, Lorentz-covariance, cross-chain operational form) --- that serves as a unified substrate mechanism for non-Markovian behavior across three physical regimes:

\begin{itemize}[noitemsep]
\item Soft-matter Maxwell viscoelastic relaxation ($\tau_{V5}^\mathrm{soft}$ identified with empirically-measured Maxwell relaxation time).
\item Hawking-spectrum cutoff ($\tau_{V5}^\mathrm{BH} = \ell_P/c$ chosen to match the substrate scale).
\item Entanglement-bandwidth modulation ($\tau_{V5}^\mathrm{ent}$ system-specific, inherited).
\end{itemize}

The three regimes have different correlation times --- \textbf{all inherited parameters, not substrate predictions}. The substrate-level structural features (kernel form, memory, support, gauge-covariance) are the same in all three. This structural unification of formalism --- a single substrate kernel rule-type underlying all three --- is the wedge claim.

What the paper does \emph{not} claim:

\begin{itemize}[noitemsep]
\item The primitives are not derived; they are postulated.
\item $\tau_{V5}$ values are inherited from value-layer empirical content, not substrate predictions.
\item V5 is not claimed as necessary for non-Markovian dynamics in some absolute sense; standard regime-specific frameworks supply alternative descriptions. V5's distinctive content is the \emph{substrate-level unification of formalism} across regimes.
\item V5 does not replace standard continuum cross-correlation physics --- it provides a unified substrate mechanism that, through DCGT, recovers standard continuum laws in each regime.
\item The Hawking cutoff scale $\omega_c \sim c/\ell_P$ matches the physical Planck scale \textbf{because} the substrate framework identifies $\ell_\mathrm{ED}$ with $\ell_P$ (Paper\_027 \S5.3); this is empirical matching, not substrate derivation.
\end{itemize}

V5 is offered for empirical engagement on these terms: a substrate primitive whose cross-regime structural-formalism unification is the empirical case, and whose specific numerical values are inherited from value-layer empirical content.

\textbf{Series context.} First Wave-2 Generative Paper. Wave-2 continues with Paper \#21 (Memory-Kernel Cascade), Paper \#22 (Lindblad Limit), Paper \#23 (Kernel Hierarchy Structural Unification).

\section*{Appendix: Cross-References and Glossary}

\subsection*{Cross-references}

\begin{itemize}[noitemsep]
\item \textbf{Position paper:} \texttt{paper\_ED\_Framework\_13\_Primitive\_Generative\_System.md}.
\item \textbf{Primitive load-bearing audit:} \texttt{PRIMITIVE\_LOAD\_BEARING\_AUDIT.md}.
\item \textbf{Paper \#1 (Participation Measure), \#3 (Inner Product), \#5 (Gauge T17).}
\item \textbf{Paper \#18 (V1 N1), Paper\_093 (T18):} kernel inheritance.
\item \textbf{Paper\_027 (Newton's $G$):} empirical identification $\ell_\mathrm{ED} = \ell_P$ that the \S5.2 Hawking matching uses.
\item \textbf{Arc-D (DCGT closure):} substrate-to-continuum bridge.
\item \textbf{Arc-Hawking (H-4 closure):} Hawking-cutoff via V5 KMS.
\item \textbf{Arc-E (E-5 closure):} entanglement-bandwidth via V5.
\end{itemize}

\subsection*{Glossary}

\begin{itemize}[noitemsep]
\item \textbf{V5 (cross-chain correlation kernel rule-type).} Substrate kernel rule-type mediating correlations between distinct chains.
\item \textbf{Kernel rule-type.} Category of substrate rule-types whose participation content is correlation amplitude between chains or substrate-graph regions. Includes V1, V5, future cascade kernels.
\item \textbf{Memory time $\tau_{V5}$.} V5-characteristic temporal scale. \textbf{Regime-specific; inherited from value-layer empirical content, not substrate-predicted.}
\item \textbf{V5 envelope $F_{V5}$.} Bounded, decaying function defining V5's specific functional form. Structural form fixed; specific shape inherited.
\item \textbf{Structural unification of formalism across regimes.} V5's cross-regime substrate-level structural-feature consistency: same substrate kernel rule-type, same structural features, different inherited correlation times. The empirical case for V5.
\item \textbf{Unified substrate mechanism for non-Markovian behavior.} What V5 provides across regimes --- one substrate-level kernel structure that supplies non-Markovian effects across diverse physical contexts.
\item \textbf{DCGT (Diffusion Coarse-Graining Theorem).} Substrate-to-continuum bridge.
\item \textbf{Imaginary-time periodicity.} Substrate-level KMS condition V5 imposes in Euclidean near-horizon BH geometry; produces Hawking spectrum at $T_H = \kappa/(2\pi)$.
\item \textbf{Lieb-Robinson velocity.} DCGT continuum-limit bound on entanglement-content propagation from V5's finite-memory cross-chain correlation.
\item \textbf{ER=EPR-class straddling.} V5-mediated cross-chain correlations spanning event horizons; substrate-level mechanism for BH information-paradox structural avoidance.
\item \textbf{Substrate-scale matching.} Operational identification of a V5 regime-specific memory time with a substrate scale, conditional on prior substrate-scale empirical identifications (e.g., $\ell_\mathrm{ED} = \ell_P$).
\item \textbf{Empirical identification.} Operational identification of a V5 regime-specific memory time with an empirically-measured timescale (e.g., $\tau_{V5}^\mathrm{soft} = \tau_\mathrm{Maxwell}$).
\end{itemize}

\end{document}
